\documentclass[11pt,a4paper]{article}
\usepackage[utf8]{inputenc}
\usepackage[margin=2cm]{geometry}
\usepackage{graphicx}\setlength{\parindent}{0pt}\setlength{\parskip}{6pt}
\usepackage{lipsum}
\newcommand{\titulo}[1]{\begin{center}\textbf{#1}\end{center}}
\begin{document}
\pagestyle{empty}
\begin{center}UNIVERSIDADE\\[8cm]\textbf{TRABALHO COM ERROS}\\[8cm]CIDADE 2026\end{center}\newpage
\titulo{RESUMO}
Este resumo é curto demais. \par Palavras-chave: teste.\newpage
\titulo{SUMÁRIO}
1 INTRODUÇÃO\dotfill 3\par
2 DESENVOLVIMENTO\dotfill 3\par
REFERÊNCIAS\dotfill 9\par
\newpage
\pagestyle{plain}\setcounter{page}{3}
\textbf{1 INTRODUÇÃO}\par
\lipsum[1-3]
\begin{center}Figura 1 – Um retângulo qualquer\\ \rule{6cm}{3cm}\end{center}
\lipsum[4]\par
\textbf{2 DESENVOLVIMENTO}\par
\lipsum[5-7]
\newpage
\textbf{REFERÊNCIAS}\par
SOUZA, A. \textbf{Livro B}. São Paulo: Editora, 2020.\par
ALMEIDA, C. \textbf{Livro A}. Curitiba: Editora, 2019.\par
MOURA, D. \textbf{Livro C}. Rio de Janeiro: Editora, 2018.\par
\end{document}
